% !PRINTED 421
\weberpage{421}
%<<CONT
in which, as we shall now often do,

\weq{8}{m = -D}

has been set, so that $m$ denotes a positive integer.

This function, which is very important for what follows, is called the \gesperrt{class function}. The roots of the equation

\weq{9}{H_m(u) = 0}

are the class invariants of the discriminant $D = -m$, and this equation is therefore called the \gesperrt{class equation}.

The degree $h$ of the class equation is equal to the number of primitive classes of the quadratic forms of the discriminant $D$. If we now consider the quotients (6), then from all this there results:

V. \gesperrt{The function $F_n(u, u)$ is divisible by $[H_m(u)]^k$, and the quotient is relatively prime to $H_m(u)$.}

Accordingly the function $F_n(u, u)$ can be decomposed into factors, and it follows, if $C$ is a constant equal to the coefficient of the highest power $u^N$ of $u$:

\weq{10}{F_n(u, u) = C\,H_m^k(u)\,H_{m'}^{k'}(u)\,H_{m''}^{k''}(u) \ \ldots,}

if the product extends over all the positive integers $m, m', m'', \ldots$ for which the equation § 114, (8)

\weq{11}{4\,n = y^2 + m\,x^2}

admits proper solutions, and $k, k', k'', \ldots$ denotes in each case the number of these solutions (with the modification indicated there in the two exceptional cases $D = -4$, $D = -3$).

The degree $N$ of the equation $F_n(u, u)$ is to be determined first.

If we choose the representatives in (5) as in § 69, then

\weq{12}{F_n[j(\tau), j(\tau)] = \Pi\left[j(\tau) - j\left(\frac{c + \partial\,\tau}{a}\right)\right].}

Here $a, \partial$ denoted all the pairs of positive integers satisfying the condition $a\,\partial = n$, % [sic] German source prints a singular verb with the plural subject.
and $c$ runs through a system of residues with respect to the modulus $a$, with the exclusion of those values which are not relatively prime to the greatest common divisor $e$ of $a$ and $\partial$, so that the number of values of $c$ belonging to a decomposition of $n$ into the two factors $a$ and $\partial$ is equal to $a\,\varphi(e){:}e$. If now $N$ is the degree of $F_n(u, u)$ and $C$ the coefficient of the highest power of $u$, then the development begins
%CONT>>
